\documentclass[12pt,a4paper]{article}
\usepackage[utf8]{inputenc}
\usepackage[T1]{fontenc}
\usepackage{amsmath,amsfonts,amssymb,amsthm}
\usepackage{mathtools}
\usepackage{microtype}
\usepackage{geometry}
\geometry{top=2.5cm, bottom=2.5cm, left=2.5cm, right=2.5cm}

\usepackage{hyperref}
\hypersetup{
    colorlinks=true, linkcolor=blue, citecolor=blue, urlcolor=blue,
    pdftitle={Weierstrass Gamma Function by Primes},
    pdfauthor={Vladislav Panin}
}

\title{Weierstrass Gamma Function by Primes}
\author{Vladislav Panin}
\date{November 10, 2026}

\begin{document}

\maketitle

\begin{abstract}
\noindent
We introduce a Weierstrass-type gamma function defined by a product over the primes,
derive an integral representation involving the prime-counting function $\pi(x)$,
and formulate a conjecture on the location of zeros in the strip $0 < \operatorname{Re}(z) < 1$.
Numerical approximations of the first zeros on the line $\operatorname{Re}(z) = \frac{1}{2}$ are presented.
\end{abstract}

\begin{center}
{\large\bfseries Definition and Integral Representation}
\end{center}

The Weierstrass gamma function by primes:
\begin{equation*}
    \frac{1}{G_p(z)} = z \prod_{p}^{\infty} \left( 1 + \frac{z}{p} \right) e^{-z/p}
\end{equation*}

\begin{equation*}
    \frac{1}{G_p(z)} = z \exp\left( - \int_{2}^{\infty} \pi(x) \frac{z^2}{x^2(x+z)} \, dx \right)
\end{equation*}

\textit{Remark:} The integral converges absolutely for $z \in \mathbb{C} \setminus (-\infty, -2]$.

\begin{equation*}
    \Phi(z) = \frac{G_p\left(\frac{1 - z}{2}\right)}{G_p\left(\frac{z}{2}\right)}
\end{equation*}

\textbf{Conjecture:} All zeros of $f(z) = 1 - \Phi(z)$ with $0 < \operatorname{Re}(z) < 1$ lie on the line $\operatorname{Re}(z) = \frac{1}{2}$.

The zeros are located at $z = \frac{1}{2} + it$, $t \in \mathbb{R}$:
\begin{align*}
    t \in \{ &7.5593242,\; 14.6100647,\; 20.3194037,\; 25.5268164,\; 30.4450039,\; 35.1672954, \\
    &39.7447722,\; 44.2091084,\; 48.5815992,\; 52.8774028,\; 57.1077783,\; 61.2813706, \\
    &65.4049958,\; 69.4841471,\; 73.5233328,\; 77.5263103,\; \dots \}
\end{align*}

\begin{thebibliography}{9}
\bibitem{froberg1968}
Fr\"oberg, C. E. (1968). 
\newblock On the prime zeta function.
\newblock \emph{BIT Numerical Mathematics}, 8(3), 187--202.

\bibitem{apostol1976}
Apostol, T. M. (1976). 
\newblock \emph{Introduction to Analytic Number Theory}. 
\newblock Springer-Verlag.
\end{thebibliography}

\end{document}
